% P65 local phase comparison candidate; independent review pending.
\paragraph{Exact comparison of the two local phases at the principal embedding.}
Use the one-based convention of P-F60:
$w_i=(-2)^{i-1}$ for $1\le i<n$, $w_n=0$,
and $\zeta=e^{2\pi i/M}$.  Define the nonzero products
\[
\begin{split}
 S&=\prod_{i=1}^{n-1}(1-\zeta^{w_i}),&
 R&=\prod_{j=2}^{n}(1-\zeta^{w_j-1}),\\
 E&=\prod_{i=1}^{n-1}(1-\zeta^{3w_i}),&
 T&=\prod_{1\le i<j<k\le n}(1-\zeta^{w_i+w_j+w_k}).
\end{split}                                                    \tag{P-F65k}
The index of $R$ starts at two, excluding $w_1=1$ and including
$w_n=0$.  Distinctness of the $w_i$ modulo $M$ makes $R,S$ nonzero;
$3$ is coprime to $M$, making $E$ nonzero.  The exact fixed-monomial
enumeration in P-F60 shows that no three distinct weights sum to
zero modulo $M$, making $T$ nonzero.
The chain component's local coefficient is exactly
\[
                    C_n^{\mathrm{chain}}(\zeta)
                         =\frac{R^3E^2}{S^3T^6}.       \tag{P-F65l}
Here is the full product cancellation.  Put
\[
 A=\prod_{i\ne j}(1-\zeta^{w_i-w_j}),\quad
 A_j=\prod_{i\ne j}(1-\zeta^{w_i-w_j}),\quad
 D=\prod_{i\ne j:\ 2w_i+w_j\not\equiv0\ (M)}
                                      (1-\zeta^{2w_i+w_j}),
\]
where $j$ is fixed in the definition of $A_j$, so $A=\prod_jA_j$.
For $1\le i\le n-2$, the relation $2w_i=-w_{i+1}$ makes
the $i$th portion of $D$ equal to
$A_{i+1}/(1-\zeta^{3w_i})$: the zero term $j=i+1$ is excluded,
and the otherwise allowed adjoint term $j=i$ is precisely the
displayed missing denominator factor.
For $i=n-1$, the portion is $S/(1-\zeta^{w_{n-1}})$, while
for $i=n$ it is $S$.  Since $\zeta^{3w_{n-1}}=
\zeta^{w_{n-1}}$ and $A_1=R,A_n=S$, multiplication gives
$D=AS/(RE)$.
The ordered tensor weights separate into pure cubes with product
$E$, two-equal triples with product $D^3$, and all-distinct triples
with product $T^6$.  The chain local formula from P-F60 is therefore
$A^3/(ED^3T^6)$, which is exactly (P-F65l).

All phases can now be computed without an asymptotic estimate.
Let $s_w=\sum_{i=1}^{n-1}w_i$; then
$3s_w=1-(-2)^{n-1}$.
For $2\le j<n$, the two numbers $w_j-1$ and $w_j$ have the
same nonzero sign and both lie strictly between $-M$ and $M$.
Their sine factors thus have the same sign.  The last factor
of $R$ is indexed by $w_n-1=-1$, whereas the first factor of
$S$ is indexed by $w_1=1$.  Since the two products have the
same length and the sums of their exponents differ by $-n$,
\[
\begin{split}
 \operatorname{phase}(R/S)&=-e^{-\pi i n/M},\\
 \operatorname{phase}(E^2)&=(-1)^{n-1}e^{2\pi i/M},\\
 \operatorname{phase}(T^6)&=(-1)^{\binom n3}
                              e^{2\pi i\binom{n-1}{2}/M}.
\end{split}                                                    \tag{P-F65m}
The last two identities follow directly from
$1-e^{2\pi ix/M}=-2i e^{\pi ix/M}\sin(\pi x/M)$:
the even powers of the sine factors are positive,
the sum of all distinct-triple exponents is
$\binom{n-1}{2}s_w$, and the integer multiples of $2\pi i$
from the displayed identity for $3s_w$ disappear.
Substitution in (P-F65l), using
$2\binom{n-1}{2}=n^2-3n+2$, gives
\[
 \operatorname{phase}(C_n^{\mathrm{chain}}(\zeta))
                =(-1)^{n+\binom n3}e^{-\pi i n^2/M}
                =\sigma_n e^{-\pi i n^2/M},            \tag{P-F65n}
\]
where $\sigma_n=-1$ exactly for $n\equiv1\pmod4$.
For the last equality, the parity of $\binom n3$ is periodic
modulo four, since
$\binom{n+4}{3}-\binom n3=4\binom n2+6n+4$ is even;
its values in residues $0,1,2,3$ are $0,0,0,1$.
This is exactly the phase in (P-F65i).  Consequently
\[
        \frac{\mathcal C_n(e^{2\pi i/M})}
             {C_n^{\mathrm{chain}}(e^{2\pi i/M})}>0.    \tag{P-F65o}
\]
The two explicitly computed local terms therefore do not cancel
each other at this embedding.  This does not control the other
fixed components, does not establish a global pole, and makes no
sign claim at other primitive embeddings.
